\section{Mitigazioni e comportamenti corretti}
\label{sec:mitigazioni}

Le misure descritte di seguito rispondono alle debolezze analizzate nelle sezioni precedenti: per ciascuna indichiamo chi deve adottarla e in quali casi non basta. Alcune sono state introdotte o raccomandate dagli stessi progetti dopo una serie di attacchi.

\subsection{Ridurre il rischio prima dell'incidente}
\label{sec:ridurre-rischio-prima}

\subsubsection{Controllo degli account e delle modifiche}
\label{sec:controllo-account}

Nel caso npm un solo account di pubblicazione, raggirato con una mail di phishing, è stato sufficiente a distribuire versioni malevole (sez. 3.1). Una prima misura è l'uso di una chiave di sicurezza (WebAuthn), che la documentazione npm descrive come metodo di 2FA~\cite{npm-2fa-docs-2026}. GitHub ha annunciato di voler sostituire il 2FA a codice temporaneo con uno basato su FIDO e raccomanda ai manutentori di usare WebAuthn [blog GitHub, 22/9/2025]. La misura spetta al manutentore. Il limite è che protegge l'account, ma non le situazioni in cui la fiducia è abusata senza rubare un account, come nell'AUR, dove nuovi account hanno adottato pacchetti orfani.

Nel caso AUR l'abuso passa da adozioni e aggiornamenti malevoli di pacchetti, e Arch raccomanda di rivedere ogni modifica a \texttt{PKGBUILD} e agli script di installazione quando si aggiorna~\cite{arch-aur-incident-2026}. La misura spetta all'utente che installa. Sono disponibili alcuni strumenti per analizzare il contenuto di \texttt{PKGBUILD}, ma non sostituiscono una revisione manuale attenta~\cite{archwiki-aur}: non tutti la fanno, soprattutto per gli aggiornamenti, e all'attaccante basta riuscirci una volta.

\subsubsection{Attendere prima di adottare nuove release}
\label{sec:attendere-release}

Secondo l'advisory ufficiale, npm ha rimosso la versione compromessa nel corso della giornata dell'8 settembre 2025~\cite{debug2025}; quindi un'attesa di almeno 24 ore avrebbe superato quella finestra. pnpm permette di impostare un'età minima delle versioni, pari a un giorno per impostazione predefinita~\cite{pnpm-supply-chain}, e Dependabot attende almeno tre giorni prima di proporre un aggiornamento, senza ritardare quelli di sicurezza~\cite{github-dependabot-cooldown}. La misura spetta a chi sviluppa il progetto. Il limite è che il ritardo riduce l'esposizione ma non la elimina: chi usa una cache o un mirror può ricevere ancora la versione compromessa~\cite{debug2025}, e la stessa documentazione di pnpm parla di probabilità, non di garanzia~\cite{pnpm-supply-chain}.

\subsection{Reagire e continuare a controllare}
\label{sec:reagire-controllare}

\subsubsection{Risposta proporzionata all'esposizione}
\label{sec:risposta-proporzionata}

Nel caso npm il registro ha rimosso la versione malevola e il proprietario del pacchetto ha pubblicato nuove versioni per svuotare le cache dei registri privati; chi gestisce mirror deve eliminare le versioni compromesse~\cite{debug2025}. Gli sviluppatori che usavano la libreria nel browser devono aggiornare all'ultima patch (ultima versione pubblicata dal progetto), rimuovere \texttt{node\_modules}, svuotare la cache del gestore di pacchetti e ricostruire i bundle~\cite{debug2025}. Il limite è che togliere il pacchetto dal registro non ripulisce i bundle già costruiti~\cite{debug2025}.

Nel caso AUR gli interventi documentati riguardano il repository: ripristino della revisione precedente e sospensione dell'account, oltre alle restrizioni temporanee decise da Arch~\cite{aur-mail-gnome-randr,arch-aur-incident-2026}. Per chi avesse già installato la versione malevola, le fonti consultate non indicano come ripristinare il sistema; poiché lo script si attiva all'installazione o all'aggiornamento, la prima verifica è controllare se il pacchetto è stato installato o aggiornato dopo l'adozione.

\subsubsection{Controllo degli avvisi e dei cambiamenti}
\label{sec:controllo-avvisi}

Sia l'advisory di debug sia l'avviso di Arch indicano azioni concrete agli utenti, quindi seguire gli avvisi ufficiali dei progetti usati è parte della difesa~\cite{debug2025,arch-aur-incident-2026}. Gli utenti AUR sono invitati a segnalare alla mailing list \texttt{aur-general} i commit sospetti sui pacchetti che usano~\cite{arch-aur-incident-2026}. Poiché l'avviso parla di adozioni malevole, un cambio di manutentore è un'anomalia da verificare prima di aggiornare.
